% R. Brown's translations of two letters to L. Breen, with his note and comments (cote 134-2): archive cote 134-2, site page 122 (PDF page 123).
% translation of the letter of 17/2/1975. Draft compared: none.
% AI-transcribed by Claude Opus 5.5 (claude-opus-5-5, adaptive thinking, effort high), Message Batches API, batch brown_full (msgbatch_01Md3i5yK1PxkpJn5GHfEExE), 2026-11-13; not yet checked by eye.
\scanpage{134-2/122}
\uline{Translation of a letter of A. Grothendieck to L. Breen 17/2/1975}

Dear Larry,

Here is an afterthought to ``une lettre-fleuve'' on the yoga of homotopy.
As you doubtless know, to a topos $X$ one associates canonically a pro-simplicial
set, and so in a convenient sense a ``pro-homotopy type''. When $X$ is ``locally
homotopically trivial'', the associated pro-object is essentially constant as
a pro-object in the homotopy category, and so $X$ defines, in the usual homotopy
category, an object which is the ``homotopy type''. Similarly, if $X$ is
``locally homotopically trivial in $\dim \leq n$'' , it defines an ordinary homotopy
type, but ``truncated in $\dim \leq n$'' - this is a familiar construction for $n = 0$
or $1$ , even among those like me who know hardly any homotopy theory!

These constructions are functorial in $X$ . Moreover, if $f : X \to Y$ is
a morphism of toposes, Artin-Mazur have given a \textit{cohomological} condition which \margin{topoi}
is necessary and sufficient for $f$ to be a ``homotopy equivalence in $\dim \leq n$'':
it is that $H^{i}(Y, F) \xrightarrow{\cong} H^{i}(X, f^{*}(F))$ for $i \leq n$ , and all \textit{locally constant}
sheaves of groups $F$ on $Y$ , allowing for $i \leq 1$ that $F$ be non-commutative.
This criterion, in terms of ``locally constant'' n-gerbes $F$ on $Y$ , can be
interpreted as the condition that $F(Y) \to F(X)$ is an n-equivalence for all
such $F$ and $i \leq n$ . It is certainly true that this is equivalent to the
following criterion:

(A) For every ``locally constant'' n-champ $F$ on $Y$ , the n-functor \margin{\uline{Stack}}
$F(Y) \to f^{*}(F)(X)$ is an n-equivalence,

or again

(B) The n-functor $F \to f^{*}(F)$ which sends the n-category of locally
constant (n-1)-champs in $Y$ to that of locally constant (n-1)-champs on
$X$ , is an n-equivalence.

In other terms, the constructions on a topos $X$ which one can make in
terms of (n-1)-champs which are \textit{locally} constant, depend only on its ``n-%
truncated pro-homotopy type'', and define it. In the case where $X$ is locally
homotopically trivial in $\dim \leq n$ , and so defines a n-truncated ordinary
homotopy type, one can interpret these last as an n-groupoid $C_{n}$ , (defined
up to n-equivalence). In terms of these

(C) The (n-1)-champs on $X$ should be able to be identified with the
n-functors from the category $C_{n}$ in the n-category ((n-1)-Cat) of all
(n-1)-categories.

In the case $n = 1$ , this is nothing other than the Poincaré theory of
the classification of coverings of $X$ in terms of the``fundamental groupoid''
$C$ , of $X$ . By extension, $C_{n}$ merits the name \textit{fundamental n-groupoid} of $X$ ,
which I propose to write $\Pi_{n}(X)$ . Knowledge of this includes knowledge of
the $\pi_{i}(X)$ $(0 \leq i \leq n)$ and the Postnikoff invariants of all orders up to
$H^{n+1}(\Pi_{n-1}(X), \pi_{n})$ .

In the case of an arbitrary topos, not necessarily locally homotopically
trivial in $\dim \leq n$ , one hopes to be able to interpret the (n-1)-champs which
are locally constant on $X$ in terms of a $\Pi_{n}(X)$ which will be a pro-n-groupoid.
This has been done, more or less, for $n = 1$ (at least for connected $X$) ;
the case where $X$ is the étale topos of a scheme is treated extensively in
SGA3, in relation to the classification of tori on an arbitrary base.

\mspage{123}
